% ROME student sheet (focus plan F14b, 20 Nov 2026). DRAFT. At most six pages (check_procedure.py C7).
% The short sheet a student follows at the robot: one card per mode, in run order. Everything else (setup, sign-offs,
% lab_settings.m, troubleshooting, the run plan, provenance) is in the TA manual, ROME_TA_manual.tex.
% Packages are limited to the Plan 3 set. Run check_procedure.py (python -I -X utf8) after every edit.
% Provenance rules (as in the TA manual):
%   * every number in the text is on number_allowlist.tsv with a file:line citation, or sits inside \todo{...};
%   * \src{path:lines} cites the code line behind a statement. This sheet is built with \showsrcfalse, so the
%     citations are checked by check_procedure.py (C4, citation_lock.tsv) but NOT printed (F13 finding 22);
%   * review documents are cited on comment lines of their own that start with "review-src:".
% The placeholders for the expected results (table and figure, focus plan F14c) are boxes of about the size the
% generated table and figure will take, so the page count measured today holds once F14c fills them.
\documentclass[10pt]{article}
\usepackage{graphicx}
\usepackage{amsmath}
\usepackage{siunitx}
\usepackage{enumitem}
\usepackage{booktabs}
\usepackage{tikz}
\usepackage{hyperref}

\setlength{\textwidth}{7in}
\setlength{\oddsidemargin}{-0.25in}
\setlength{\evensidemargin}{-0.25in}
\setlength{\textheight}{9.4in}
\setlength{\topmargin}{-0.7in}
\setlength{\parindent}{0pt}
\setlength{\parskip}{3pt}
\setlength{\emergencystretch}{3em}
\hypersetup{colorlinks=true, linkcolor=blue, urlcolor=blue}
% Let the Mode 1 table and figure share a page with text (no page of floats alone).
\renewcommand{\topfraction}{0.9}
\renewcommand{\bottomfraction}{0.9}
\renewcommand{\textfraction}{0.05}

% ---------------------------------------------------------------- macros (same as the TA manual)
\newcommand{\todo}[1]{\textcolor{red}{\textbf{[#1]}}}
\ExplSyntaxOn
\cs_new_protected:Npn \romecode #1
  {
    \str_map_inline:nn {#1}
      {
        ##1
        \str_if_in:nnT { /_;:,. } {##1} { \penalty 50 \relax }
      }
  }
\ExplSyntaxOff
\newcommand{\cmd}[1]{\texttt{\romecode{#1}}}
\newcommand{\out}[1]{\texttt{\romecode{#1}}}
\newif\ifshowsrc
\showsrcfalse
% Hidden on this sheet: \unskip removes the space before it, so no stray space is printed before a full stop.
\newcommand{\src}[1]{\ifshowsrc{\scriptsize\,[\texttt{\romecode{#1}}]}\else\unskip\fi}
\newcommand{\expect}{\par\textit{You see:}\ }
% A framed box the width of the text.
\newcommand{\framed}[1]{\par\smallskip\noindent\fbox{\parbox{\dimexpr\linewidth-2\fboxsep-2\fboxrule\relax}{#1}}\par\smallskip}
% The TA's sign-off box at the top of a card.
\newcommand{\prereq}[1]{\framed{\textbf{TA signs before you start this card.} #1}}
% What success looks like, and what to do otherwise, at the end of a card.
\newcommand{\success}[2]{\framed{\textbf{Success:} #1\\\textbf{Otherwise:} #2}}
% A placeholder of a fixed size for a generated table or figure that does not exist yet (F14c).
% #3 is the \todo{BLOCKED: ...} itself, so check_procedure.py C2 sees it in the \IfFileExists fallback.
\newcommand{\placeholder}[3]{\fbox{\parbox[c][#2][c]{#1}{\centering #3}}}
% Numbered steps of a card; the numbering restarts on every card.
\newlist{steps}{enumerate}{1}
\setlist[steps]{label=\arabic*., leftmargin=*, itemsep=2pt, topsep=2pt}

\title{\vspace{-1.2cm}ROME robot: student sheet}
\author{Spacecraft simulation laboratory}
\date{DRAFT, 20 November 2026}

\begin{document}
\frenchspacing
\maketitle
\vspace{-0.6cm}

% What a robot card says about the verdict line (the same on modes 2, 3 and 4).
\newcommand{\verdictstep}[1]{\item \textbf{Read the result.} Within about \SI{10}{\second} of the stop MATLAB prints a
  line that begins \out{RUN VERDICT mode #1:} \src{rome_run_end.m:14; run_verdict_core.m:161} On \out{FAIL} or
  \out{INCOMPLETE}: start nothing else, show the line to the TA, and go on only when the TA says so. No such line within
  about \SI{10}{\second}: call the TA.}
% The calibration step, written out on every robot card (no pointer to another card).
% Success is judged by the park-record line only (clarity audit A2): the "successfully completed" line is not used.
\newcommand{\calibbody}{The arm moves now, on every card: each joint drives into its limit switch. Clear zone around the
  arm empty; your partner's hand on the power cut-off; only the power cut stops a calibration. Type
  \cmd{calibrateROMEArm;} and wait; it can take up to three minutes \src{calibrateROMEArm.m:42; calibrateROMEArm.m:11; calibrateROMEArm.m:133}.
  \textbf{Success:} near the end MATLAB prints \out{Park record written (entry <n>)} \src{calibrateROMEArm.m:216}. If it
  prints \out{NO park record written} \src{calibrateROMEArm.m:218} or \out{DO NOT run a robot mode}
  \src{calibrateROMEArm.m:192}, or the arm does not look like the text under \out{Expected posture}
  \src{calibrateROMEArm.m:189-190}, stop and call the TA. Never move a joint yourself (\cmd{jog_joint}, \cmd{jogROMEArm}):
  that is a TA action \src{calibrateROMEArm.m:193}.}
% Placement with the power off (author, 20 Nov 2026: "Power off to place").
\newcommand{\placebody}[1]{\textbf{Place the robot with its power off.} Your partner turns the robot's power off with
  the switch the TA showed you; the motors then have no power. Push it by hand #1 Then your partner turns the power on.
  Before you calibrate, check that the lab PC is still on the robot's Wi-Fi (the TA shows you where) \src{docs/lab/PART_F_BENCH_CHECKLIST.md:66}.
  \todo{NOT VERIFIED: whether the power switch also turns off the robot's Wi-Fi, Part F item 25c}}

% =====================================================================================
\section*{How to use this sheet}
Do the cards in order: Mode 1, then Mode 2, then Mode 3, then Mode 4. Each card is one run, and its steps are in the
order you do them: set the mode, place the robot, calibrate the arm, Run, read the result, save. Start a card only
after the TA has signed its box. The TA has a separate manual with everything else.
% review-src: docs/superpowers/plans/2026-11-19-case1-three-modes.md:87

Type every command exactly as printed, including the semicolon at the end where there is one. Two kinds of lines are expected and can be
ignored: the warning \out{The orbit exceeds a 2m x 2m area!} \src{define_constants.m:335}, and at the Run the lines
that begin \out{*** ARM FROZEN} or \out{*** ARM WILL MOVE} \src{arm_freeze_banner.m:144; arm_freeze_banner.m:147}.

\textbf{You do not edit files.} You set the mode by typing \cmd{rome_mode(n);} (or click the \textbf{Mode n} button at
the top of the model) \src{tools/wire_run_mode_buttons.m:38}. Only if the TA asks: type \cmd{edit define_constants},
change \cmd{speed_factor} or \cmd{size_factor} \src{define_constants.m:8-11}, save (Ctrl+S), and start again at the
Mode 1 card.

\framed{\textbf{Red text, an error window, or an orange warning this sheet does not tell you to expect: stop and call
the TA.} Do not try to work around it. The common ones: \out{Run refused.} or \out{PRERUN CHECK mode <n>: REFUSED}
(the robot is not ready) \src{prerun_check.m:101; prerun_check.m:103}; \out{Failed to establish TCP connection} (no
link to the robot) \src{calibrateROMEArm.m:44}; \out{Calibration timeout} \src{calibrateROMEArm.m:113};
\out{READY status not confirmed by Teensy.} \src{calibrateROMEArm.m:202}; \out{stop it before changing the run mode}
\src{rome_mode.m:23}.}
% review-src: TRACKER.md section 0

\subsection*{Safety}
\begin{itemize}[leftmargin=*, itemsep=1pt, topsep=1pt]
  \item Work in pairs. During every motion your partner keeps a hand on the robot's power cut-off and you are at the
  keyboard \src{docs/lab/PART_F_BENCH_CHECKLIST.md:18}.
  \todo{BLOCKED: where the power cut-off is, Part F item 25c}
  \item Keep the clear zone around the arm empty while the arm has power
  \src{docs/lab/PART_F_BENCH_CHECKLIST.md:20}. Stay clear below the arm when the power is cut: whether it holds or
  falls without power is not known \src{docs/lab/PART_F_BENCH_CHECKLIST.md:21}.
  \todo{BLOCKED: clear-zone radius and arm behaviour on power loss, Part F items 24d and 25c}
  \item \textbf{To stop:} (a) press \textbf{Stop} in Simulink; that sends the stop command to the robot
  \src{ROMECommand.m:197-200}. (b) If anything still moves, cut the power.
  \item \textbf{During arm calibration only the power cut stops the arm.} Nothing you type stops it
  \src{UK-Arduino/ROME_Teensy_Code/ROME_Teensy_Code.ino:160-163}.
  \item \textbf{Never move a joint yourself} (\cmd{jog_joint}, \cmd{jogROMEArm}). Moving a joint is a TA action, also
  when MATLAB prints a joint problem \src{calibrateROMEArm.m:193}. Call the TA.
  \item After any stop, start nothing (no run, no calibration, no joint move) until the TA says so: the arm moves
  again on the next joint command \src{UK-Arduino/ROME_Teensy_Code/ROME_Teensy_Code.ino:181-188}. After a power cut
  the arm must be calibrated again \src{UK-Arduino/ROME_Teensy_Code/ROME_Teensy_Code.ino:52}. Tell the TA about every
  stop, with the time and what you saw.
\end{itemize}
% review-src: TRACKER.md section 0, row 19 Nov "Arm does not stay stopped after STOP_ALL"; TRACKER.md section 0, row 19 Nov "Nothing but the power cut"

\subsection*{Words used on this sheet}
\begin{description}[leftmargin=0pt, itemsep=0pt, topsep=1pt, font=\normalfont\bfseries]
  \item[Base, arm, end effector, orbit.] ROME is a four-wheeled base that carries a six-joint arm (joints J1 to J6).
  The end effector is the tool at the end of the arm. The model commands only the end effector: it follows an orbit,
  the path of a spacecraft around a target scaled down to the floor. The controller shares that motion between the
  base and the arm \src{lab_settings.m:23-28}.
  \item[Run mode.] Which of four ways the model runs (the cards below).
  \item[Model.] The Simulink model \cmd{ROME_9DOF.slx}.
  \item[Fake camera.] With the camera off, the model sends its own predicted base pose through the camera's software
  path. Nothing measures where the real base is \src{run_mode_table.m:7}.
  \item[Start card.] The block headed \out{PLACE THE ROBOT LIKE THIS} that MATLAB prints after a mode is set: the base
  position $x$, $y$ in metres, the yaw (the heading of the base) in degrees, and the six arm angles
  \src{define_constants.m:318-324}.
  \item[Calibration, limit switch.] \cmd{calibrateROMEArm} drives every joint to its limit switch (a switch at the end
  of the joint's travel, which tells the arm where it is) and then parks the arm at the start posture
  \src{calibrateROMEArm.m:1}.
  \item[Verdict line.] After a robot run MATLAB prints a line that begins \out{RUN VERDICT mode}, then the mode and
  \out{PASS}, \out{FAIL} or \out{INCOMPLETE}, and one line per item below it \src{run_verdict_core.m:149-162}.
\end{description}

\subsection*{Start of the session}
\begin{steps}
  \item Start MATLAB R2026a (the TA shows you which icon). Type \cmd{cd '<path the TA wrote on the board>'} and press
  Enter.
  \todo{BLOCKED: path of the model folder on the lab PC}
% review-src: TRACKER.md section 2, 19 Nov "Lab PC, R2026a"
\end{steps}
% The model is opened in the Run step of the robot cards, after calibrateROMEArm has put Optitrack/ on the path
% (calibrateROMEArm.m:19); opened by hand before that, MATLAB warned that the System object MotiveData cannot be found
% (F13 walkthrough finding 14). TA manual, sec:cards.

% =====================================================================================
\section*{Mode 1: simulation only}
\textit{Nothing is sent to the robot. Only MATLAB runs.} \src{run_mode_table.m:30}
\prereq{Nothing to sign for this card.}
\begin{steps}
  \item \textbf{Set the mode.} Type \cmd{rome_mode(1);}
  \expect last, \out{rome_mode: MODE 1: Simulation only (end-effector dynamics sim, pose data from sim)}
  \src{rome_mode.m:51; run_mode_table.m:30}.
  \item \textbf{Run.} Type \cmd{check_9dof;} It simulates the run and checks it \src{check_9dof.m:1-7}.
  \item \textbf{Read the result.} Near the end MATLAB prints \out{RESULT: PASS} between two lines of \out{=}
  \src{check_9dof.m:254-260}. Write down the number in the line \out{base travels to ... m from the origin} for the TA
  \src{check_9dof.m:251}.
  \item \textbf{Look at the path.} Type \cmd{plot_9dof;} It opens many figures; look only at the one titled
  \out{Desired vs. Tracked Position}, with a blue circle labelled \out{Focus}: the tracked path lies on the desired one
  \src{plotting.m:63-64; plotting.m:76}. With the default \cmd{speed_factor} and \cmd{size_factor}, compare it with
  Figure~\ref{fig:expected1} and the numbers \cmd{check_9dof} printed with Table~\ref{tab:expected}; while they are
  empty boxes, skip the comparison.
  \item \textbf{Save.} Type \cmd{save_run;} It saves the run that \cmd{plot_9dof} made and prints \out{save_run:} and
  the file name \src{save_run.m:95}. Write the file name in the lab notebook.
\end{steps}
\success{\out{RESULT: PASS}, and the tracked path lies on the desired path.}{call the TA (TA manual,
``Troubleshooting'').}

% Not floats: the table and the figure stay inside the Mode 1 card (numbered by hand with \refstepcounter).
\par\medskip\noindent
\begin{minipage}[t]{0.55\linewidth}\centering
  \refstepcounter{table}\label{tab:expected}Table~\thetable: Expected results (mode 1), default settings.\\[3pt]
  \IfFileExists{generated/expected_results.tex}{\input{generated/expected_results.tex}}{\placeholder{0.95\linewidth}{2.2in}{\todo{BLOCKED: generated from the model (the paper's results table, per case)}}}
\end{minipage}\hfill
\begin{minipage}[t]{0.42\linewidth}\centering
  \refstepcounter{figure}\label{fig:expected1}Figure~\thefigure: Expected path of case 1.\\[3pt]
  \IfFileExists{figures/fig_9dof_case_1_path.pdf}{\includegraphics[width=0.95\linewidth]{figures/fig_9dof_case_1_path.pdf}}{\placeholder{0.95\linewidth}{2.2in}{\todo{BLOCKED: expected path figure of case 1, generated from the model}}}
\end{minipage}
% provenance: fig_9dof_cases.m:10-11 (file names). The other cases are in the TA manual.

% =====================================================================================
\clearpage
\section*{Mode 2: ground vehicle moving, frozen arm, camera off}
\textit{The base drives. The arm stays parked at its start posture, so the end effector is not tracked: this is a
base-path test (only the path of the base is checked).} \src{run_mode_table.m:31; tools/wire_motive_emulator.m:70}
% review-src: TRACKER.md section 2, 19 Nov "Mode 2 is the base-path test, as built"
\prereq{Full Part F: the whole bench checklist is signed \src{docs/lab/PART_F_BENCH_CHECKLIST.md:277}. The TA has shown
you the power cut-off, set \cmd{speed_factor} for a first run, and connected the lab PC to the robot's network. The
Mode 1 card passed with the same settings.}
% review-src: TRACKER.md section 2, 19 Nov "Q-P4-a answered"
\begin{steps}
  \item \textbf{Set the mode.} Type \cmd{rome_mode(2);} (or click the \textbf{Mode 2} button at the top of the model).
  \expect last, \out{rome_mode: MODE 2: ground vehicle moving, frozen arm, camera off (fake camera)}, and the start
  card above it \src{rome_mode.m:51; run_mode_table.m:31}.
  \item \textbf{Clear the floor.} Clear the area the TA marks for the run. Nobody and nothing inside it.
  \item \placebody{to the spot and heading the TA marked; close by eye is enough.}
  \item \textbf{Calibrate the arm.} \calibbody
  \item \textbf{Change nothing now.} Do not edit a file, type \cmd{rome_mode}, or click anything on the model except
  Run. If you did, calibrate again.
% review-src: TRACKER.md section 3, Q-F11-b
  \item \textbf{Run.} If the model is not open, type \cmd{open_system('ROME_9DOF')}. Your partner's hand on the power
  cut-off; tell your partner, then press \textbf{Run}. Before the robot moves MATLAB prints
  \out{PRERUN CHECK mode 2: PASS}; the words \out{UNSOURCED} and \out{Not enforced} on these lines are normal. If it
  prints \out{REFUSED}, nothing moves: call the TA \src{prerun_check.m:95-101}. Watch the robot for the whole run.
  Stop (Safety) if any part comes near the edge of the marked area, if the arm moves, or if anything looks wrong. The
  run stops by itself at the end of the path and sends the stop command \src{ROMECommand.m:197-200}.
  \verdictstep{2}
  \item \textbf{Save.} Type \cmd{save_run;} The run was already saved at its end.
  \expect \out{save_run: already saved at the end of the run:} and the file name \src{save_run.m:119}. Write the
  file name in the lab notebook.
\end{steps}
\success{\out{RUN VERDICT mode 2: PASS}, the base drove its path, and the arm did not move.}{\out{FAIL} or
\out{INCOMPLETE}: show the line to the TA (TA manual, ``Reading the verdict line'').}

% =====================================================================================
\clearpage
\section*{Mode 3: ground vehicle moving, moving arm, camera off}
\textit{The base and the arm move and the tool follows the orbit. The fake camera gives the base pose; the controller
reads the real arm angles.} \src{run_mode_table.m:32; tools/wire_motive_emulator.m:70}
\prereq{Full Part F signed, including its joint-map test (it checks that the arm's joints J2, J5 and J6 turn the way
the model believes) \src{docs/lab/PART_F_BENCH_CHECKLIST.md:180; docs/lab/PART_F_BENCH_CHECKLIST.md:192}. The lab PC is
on the robot's network. The TA accepted your Mode 2 run.}
\begin{steps}
  \item \textbf{Set the mode.} Type \cmd{rome_mode(3);} (or click the \textbf{Mode 3} button at the top of the model).
  \expect last, \out{rome_mode: MODE 3: ground vehicle moving, moving arm, camera off (fake camera)}
  \src{rome_mode.m:51; run_mode_table.m:32}.
  \item \textbf{Clear the floor} the TA marks. Nobody and nothing inside it.
  \item \placebody{to the spot and heading the TA marked; close by eye is enough.}
  \item \textbf{Calibrate the arm.} \calibbody
  \item \textbf{Change nothing now.} Do not edit a file, type \cmd{rome_mode}, or click anything on the model except
  Run. If you did, calibrate again.
  \item \textbf{Run.} If the model is not open, type \cmd{open_system('ROME_9DOF')}. Your partner's hand on the power
  cut-off; tell your partner, then press \textbf{Run}. Before the robot moves MATLAB prints
  \out{PRERUN CHECK mode 3: PASS}; the words \out{UNSOURCED} and \out{Not enforced} on these lines are normal. If it
  prints \out{REFUSED}, nothing moves: call the TA \src{prerun_check.m:95-101}. The robot may then wait a few seconds
  before it starts \src{allowed_ranges.m:148}. Watch for the whole run. Stop (Safety) if any part comes near the edge
  of the marked area, if the tool heads toward the base or the floor, or if anything looks wrong.
  \verdictstep{3}
  \item \textbf{Save.} Type \cmd{save_run;}
  \expect \out{save_run: already saved at the end of the run:} and the file name \src{save_run.m:119}. Write the
  file name in the lab notebook.
\end{steps}
\success{\out{RUN VERDICT mode 3: PASS (1 item(s) not judged: UNSOURCED)} \src{run_verdict_core.m:160-161} (the
UNSOURCED part is expected today), and the tool followed the orbit.}{\out{FAIL} or \out{INCOMPLETE}: show the line to
the TA (TA manual, ``Reading the verdict line'').}

% =====================================================================================
\clearpage
\section*{Mode 4: ground vehicle moving, moving arm, camera on}
\textit{Mode 3 with the real OptiTrack cameras: they see markers on the robot and report where it really is on the
floor.} \src{run_mode_table.m:33; lab_settings.m:39-43}
\prereq{Full Part F signed, and also the lab-frame checklist, Part C, with the camera calibration
\src{lab_settings.m:44-48}. The lab PC is on the robot's network and the camera stream. The TA accepted your Mode 3
run.}
\begin{steps}
  \item \textbf{Set the mode.} Type \cmd{rome_mode(4);} (or click the \textbf{Mode 4} button at the top of the model).
  \expect as the last line,\\\out{rome_mode: MODE 4: ground vehicle moving, moving arm, camera on (OptiTrack)}
  \src{rome_mode.m:51; run_mode_table.m:33}. If the camera calibration is missing, Mode 4 is refused and the previous
  mode is set back \src{rome_mode.m:36-47}: stop and call the TA.
  \item \textbf{Clear the floor} around the taped origin, as the TA marks it. Keep the camera area clear: a person in
  it can hide the markers \src{docs/lab/PART_I_RUN_PLAN.md:123}.
  \item \placebody{onto the taped start mark of case 1, its heading on the heading line. If the start card and the label
  of the mark disagree, or a mark is missing or loose, call the TA \src{docs/lab/PART_C_LAB_VISIT.md:192}.}
  \todo{BLOCKED: placement tolerance, and which point of the base goes on the mark (Part C C9.5)}
  \item \textbf{Calibrate the arm.} \calibbody
  \item \textbf{Change nothing now.} Do not edit a file, type \cmd{rome_mode}, or click anything on the model except
  Run. If you did, calibrate again.
  \item \textbf{Run.} If the model is not open, type \cmd{open_system('ROME_9DOF')}. Your partner's hand on the power
  cut-off; tell your partner, then press \textbf{Run}. Before the robot moves MATLAB prints
  \out{PRERUN CHECK mode 4: PASS}; the words \out{UNSOURCED} and \out{Not enforced} on these lines are normal. If it
  prints \out{REFUSED}, nothing moves: call the TA \src{prerun_check.m:95-101}. The robot may then wait a few seconds
  before it starts \src{allowed_ranges.m:145; allowed_ranges.m:148}. Watch for the whole run. Stop (Safety) if the
  cameras lose the robot, if any part comes near the edge of the marked area, if the tool heads toward the base or the
  floor, or if anything looks wrong \src{docs/lab/PART_I_RUN_PLAN.md:123}.
  \verdictstep{4}
  \item \textbf{Save.} Type \cmd{save_run;}
  \expect \out{save_run: already saved at the end of the run:} and the file name \src{save_run.m:119}. Write the
  file name in the lab notebook.
\end{steps}
\success{\out{RUN VERDICT mode 4: PASS (1 item(s) not judged: UNSOURCED)} \src{run_verdict_core.m:160-161} (the
UNSOURCED part is expected today), and the tool followed the orbit.}{\out{FAIL} or \out{INCOMPLETE}: show the line to
the TA (TA manual, ``Reading the verdict line'').}

% =====================================================================================
\section*{End of the session}
\begin{steps}
  \item Type \cmd{rome_mode(1);} so that the next person starts in simulation \src{run_mode_table.m:30}.
  \item Write in the lab notebook the case, the modes, the settings you changed, the run file names and anything
  unusual. Hand it to the TA.
  \item Tell the TA you are done. The TA turns the robot off.
\end{steps}

\end{document}
